\initial{A} fenntartható pénzügyek kutatási iránya mind akadémiai szempontból, mind törvényhozói szempontból egyre meghatározóbb szerepet tölt be. Az 1970-es évek eleje óta készültek kezdetleges kutatások ebben az irányban, viszont az elmúlt 5 évben vált az egyik legdinamikusabban fejlődő részévé a kutatásoknak . A szakirodalmak egyik fő jellemzője az, hogy egyszerre több, egymást részlegesen fedő irányzat is megjelent hasonló fókusszal. Ez köszönhető a lefedett tématerületek gyors változásának, valamint a kutatók saját taxonómiájának\footnote{A taxonómia az adott kutatási terület egyedeinek a csoportosítási szabályait jelenti, jelen esetben a fenntarthatónak tekinthető tevékenységek felderítését és osztályozását.} is. \parencite{sarvari_fenntarthato_2026}

A fenntartható pénzügyek irányzat a fejlődése alatt sok állomáson ment át, amik közül az egyik legújabb - az általunk is kutatott - biodiverzitási pénzügyek (biodiversity finance) terulet. \parencite{karolyi_biodiversity_2023}

\subsection{Kutatási irányok és fókuszok}

Sok különböző fogalmon keresztül próbálták megragadni a kutatók a területet. Kezdetben a környezeti pénzügyek (environmental finance) fogalma jelent meg, majd ezt váltotta a karbonpénzügyek (carbon finance). Ez jól mutatja, hogy a kezdetleges környzezeti pénzügyek területéhez képest,  már itt elkezdett kialakulni a területen aprózódás. A karbonpénzügyek már kifejezetten a GHG (üvegházhatású gázok) által okozott szennyezéseket és ezeknek pénzügyi hatásait vizsgálja.

A következő lépés 2006 után az úgynevezett klímapénzügyek terület lett, a karbonpénzügyek területe pedig kezdett visszaszorulni. A klímapénzügyek irányzata a kibocsátáscsökkentés mellett a klímaváltozás jelentette pénzügyi kockázatok csökkentését is számszerűsíti. A gyakorlatban a klímakockázatok két csoportra bonthatók: a fizikai kockázatra (physical climate risk), ami a vállalatok termelésére közvetlen hatással bír; és az átállási kockázatra (transition risk), ami pedig az átállási költségek miatti cash-flow változásokat méri \parencite{giglio_climate_2021}. Egy fontos alága a klímapénzügyekhez az impact investing, azaz az a befektetési mód, ami a befektetők által szeretné ösztönözni a vállalatok klímaátállását. \textcite{de_angelis_climate_2021} kutatása szerint ez egy létező hatás, ami sokkal erősebben is tud érvényesülni, amennyiben a befektetők tisztában vannak azzal, hogy a kormányzat a törvényhozásában is szigorúbban jár el. Továbbá azt is kimutatták, hogy amennyiben a "zöld befektetők" által kezelt eszközök állománya megkétszereződik, a vállalatok \textit{karbonintenzitása}\footnote{Egy vállalat karbonintenzitása a vállalati kibocsátási rátájának és az egy éves nyereségének a hányadosa. \parencite{de_angelis_climate_2021}} körülbelül 5 százalékkal csökken.

A zöld pénzügyek (green finance) és a jelengleg is használt fenntartható pénzügyek (sustainable finance) kifejezések 2015 óta váltak egyre gyakoribbá. Az előző irányokhoz képest a karbonkibocsátásnál és éghajlatváltozásnál tágabb témakörben a pénzügyi rendszer egészét kutatják fenntartható szempontok alapján, mint például az alacsony kibocsátású és erőforrás-hatékony gazdaságra való átállást, környezetvédelmi projektek finanszírozását és a pénzforrások (cash-flow) \textit{"zöldítését"} \parencite{sarvari_fenntarthato_2026}. Győztesként a fenntartható pénzügyek kifejezés került ki, napjainkban a legtöbb releváns szakirodalom ezt a kategorizálást alkalmazza, jelen dolgozatot is ezzel a fogalommal kategorizáljuk

A fenntartható pénzügyek (vagy zöld pénzügyek) terület mind a kötvény-, mind a részvénypiacon is több témát vizsgált. A kötvénypiac egyik fő kérdésében a \textit{zöld kötvények (green bonds)} állnak, olyan kötvények, amik a kibocátási értékülből minősített környezeti vagy klíma vonatkozású célokat támogatnak. Fontos megemlíteni, hogy a \textit{greenwashing}\footnote{Az a jelenség, amikor egy vállalat a saját tevékenységét hamisan zöldként, vagy fenntarthatóként mutatja be.} jelenség elkerülése érdekében Európai Unióban egységes szabályozás létezik, mind az Amerikai Egyesült Államokban több különálló jogalkotásban szabályozott keret között lehet zöld kötvényeket kibocsátani. Az Európai Unióban zöld kötvényként kibocsátott papír az (EU) 2023/2631 számú rendelet alapján történik, ami opcionálisan előírja, hogy a kötvényekből származó bevételeknek 85 százalékban kell fenntartható célú\footnote{A fenntartható befektetéseket az Európai Unió az (EU) 2020/852 rendelete (EU Taxonómia rendelet) alapján jelöli ki. A rendeletben meghatározott környezeti célkitűzések között példaként szerepel az éghajlatváltozás mérséklése, az éghajlatváltozáshoz való alkalmazkodás és a biológiai sokféleség és az ökoszisztémák védelme és helyreállítása is.} befektetéseket finanszírozni. Az így kibocsátott kötvények után évente jelenteni kell a kötvény ellenértékének allokációról, majd miután ez teljesen megtörtént, a lejárat előtt közzé kell tenni egy jelentést az elért környezeti hatásokról is. A kibocsátás előtt egy külső vizsgáló fél hozzájárulására szükség van. Az USA-ban kicsocsátott zöld kötvények kibocsátásának szabályozására nincsen egy központi rendelet, erre a szerepre a legesélyesebb jelölt az ICMA Green Bond Principles rendelet, ami egy irányelvet ad a zöld kötvény kibocsátáshoz, de nem szabja meg a fenntarthatónak számító befektetéseket.

A zöld kötvénypiacon az egyik kiemelkedő kutatási irány, a kötvényken vizsgálható \textit{greenium}, azaz \textit{green bond premium}, amit a \textit{zöld kötvények} és a konvencionális kötvények közötti hozamkülönbségként értelmeznek. Több módszertannal is készültek kutatások a témában, \textcite{fandella_uncovering_2024} szerint a \textit{greenium} létezik, az értékét három különböző módszerrel is megbecsülték, viszont az eredmények különbözősége miatt nem lehet egyértelműen kimondani, hogy a zöld kötvények forgalomba hozása egyértelműbb hozamelőnyt tud biztosítani. \textcite{zerbib_is_2016} 2013 és 2017 között vizsgálta a \textit{greenium}-ot hasonló tulajdonságokkal kiadott konvencionális kötvényekhez képest. Az állapították meg, hogy a zöldkét értékelt kötvények átlagos -2 bázisponttal alacsonyabb hozamok mellett kereskednek, azonban ezt a zöld megnevezésen kívül nagyban befolyásolja a minősítő vállalat és a kötvényt kiállító vállalat típusa is.

A kötvénypiac mellett a részvénypiacon is készültek kutatások fenntarthatósági témakörökben. Az egyik legmeghatározóbb mutatócsoport az ESG pontszám (ESG Score), ami egy-egy vállalatot értékel 3 oszlop (pillar) alapján: \textit{környezeti (environment), szociális (social) és irányítási (governance)}. A vállalatok értékelését független ESG minősítők végzik (egy adott vállalatra több minősítő is ad meg értékeket), a legmeghatározóbb ESG-minősítők az MSCI, Sustainalytics, S\&P és a Refinitiv. Mivel több független minősítő cég is jelen van a piacon, az általuk adott ESG pontszámok az alkalmazott módszertanuok miatt jelentősen (akár szubjektíven) eltérhetnek egymástül, amit érdemes észben tartanunk. \textcite{berg_aggregate_2022} kutatásukban 6 prominens ESG-minősítő által becsült értékeket hasonlítottak össze, ahol azt állapították meg, hogy a pontokban az eltéréseket 58 százalékban a mérési folyamat, 38 százalékban a mért területek besorolása és 6 százalékban az alkalmazott súlyok adják. Emellett az eltéréseket tovább vizsgálva egy úgynevezett \textit{"minősítői hatást" (rater effect)} találtak, azaz egy minősítő véleménye a teljes cégről befolyásolja az általa adott ESG értéket is.

Az ESG értékek elemzésének irodalmában mindhárom oszlop alapján sok-sok különböző irányú kutatást lehet végezni, ezeknek az irányoknak a bemutatása nem célja a dolgozatunknak. A számunka releváns irányok az ESG alapú befektetések és kockázatok elemzésének irodalma. A modern ESG alapú befektetési stratégiák összegyűjtésére \textcite{amel-zadeh_why_2017} egy kérdőíves kutatásban kérdezték meg befektetési vállalatok szenior kollégáit. Leggyakrabb esetben - ügyfelek személyes kérésein és kifejezett ESG termékeken túl - akkor használnak ESG értékeket amikor valamilyen etikai szempontot is figyelembe vesznek. Tapasztalataik alapján az ESG integráció\footnote{Az ESG integráció az a befektetési folyamat, amikor az ESG értékek alapján is rangsorolják az adott részvényeket vagy kötvényeket. Leggyakrabb esetben kockázatkezelési vagy befektetői megfeleléskor alkalmazzák.} a leghatékonyabb módja a pontszámok figyelembe vételének, de emellett a negatív szűrést\footnote{A negatív szűrés az a folyamat, amikor egy meghatározott ESG érték alatt nem nem fektetnek egy vállalatba, azt kiszűrik a befektetési univerzumukból.} (negative screening) is gyakran alkalmazzák.

\textcite{giese_foundations_2019} az ESG karakterisztikák és a vállalati érték, teljesítmény közötti kapcsolatot vizsgálta. Az elemzésben három \textit{transzmissziós csatornát}\footnote{A transzmissziós csatorna az az útvonal, amik mentén egy exogén változó kifejti a hatását a vizsgált változóra.} azonosítottak, a \textit{cash-flow csatornát}, az \textit{egyedi kockázat (idiosyncratic risk) csatornát}, és az \textit{értékelési (valuation) csatornát}. Az MSCI ESG minősítési alapján vizsgálták a vállalatokat, az ESG értékeknek szignifikáns hatása volt a vállalatértékre és teljesítményre, mind a szisztematikus kockázati csatornán (alacsonyabb tőkeköltséget értek el, magasabb külső értékeket kaptak), mind pedig az egyedi kockázati csatornán (magasabb profitráta, alacsonyabb kitettség a \textit{tail kockázatokra}\footnote{A vállalati részvényértékéhez kapcsolódó kockázat, amikor a részvény értéke a historikus szórásának 3-szorosánál nagyobb árváltozást szenved el.}). A valóságban nehéz a két hatást elválasztani egymástól, mivel a magasabb ESG érték okozhat magasabb vállalati értéket, viszont ez ugyanúgy fodítva is elmondható. Mindezt figyelembe véve azt jelentették ki, hogy egy vállalat ESG karakterisztikáinak a változása hatással lehet a cég pénzügyi mutatóira.

Az ESG befeltetéssel kapcsolatos kutatások nagyrésze a befektetéseket hozam-orientált nézőpontból kutatta. \textcite{pedersen_responsible_2021} a kutatásában a befektetők egyéni hasznossági függvényét alakította át úgy, hogy a befektető által tartott portfólió magas ESG értéke pozitív hasznosságot ad. A befektető a portfólióját Markowitz-féle portfólióelmélet alapján építi, a kiegészített hasznossági függvény alapján. Azt találták, hogy a fenntartható befektetői trade-off probléma vizualizálható az ESG-hasznossági határral, ami az általa elérhető befektetési univerzumra tartalmazza az elérhető legnagyobb Sharpe-rátát. Az általuk bevezetett módszertannal egyszerre lehet az ESG információ többletértékét (mennyivel növekszik a maximálisan elérhető Sharpe-ráta az ESG információk figyelembevételével) és az ESG preferenciájának a költségét (a preferenciáji miatt mekkora Sharpe-ráta csökkenést szenved el a befektetési lehetőségeinek kiyárásával).

\subsection{Biodiverzitási kockázatok/biodiverzitási pénzügyek}

A fenntartható pénzügyek kutatási terulet jelenleg kiemelet figyelmet élvező területe a biodiverzitási pénzügyek (biodiversity finance) irány. A biodiverzitási kockázat alapjaiban más, mint a klímaváltozás kockázata. Míg a többlet széndioxid megfelelő technológiával és beruházásokkal kivonható a légkörből, ezzel szemben egy faj pusztulása végleges. Szakértői becslések szerint a világ 1000-szer nagyobb mértékben veszít el fajokat a \textit{normális kihalási rátához}\footnote{A normális kihalási ráta, vagy háttér-kihalási ráta (background extintion rate) az a természetes folyamat, ami az emberi beavatkozás nélküli fajpusztulás mértékét jelenti, ehhez viszonítva szokták értelmezni az ember vagy természet által okozott hatásokat.}  viszonyítva, és az emlősök, madarak, halak és hüllők területi előfodulása???? a legutóbbi 4 évtizedhez képest 60 százalékkal szorult vissza. \parencite{karolyi_biodiversity_2023}

Definíciója többek szerint, flammer, giglio

Relatíve friss, itt vannak ez eddigi (nem biodiv premium!) eredmenyek.

\subsection{Biodiverzitási prémium az eddigi kutatások alapján}

helena, giglio
